\documentclass[UTF8,12pt]{ctexart}

\usepackage[a4paper,margin=2.2cm]{geometry}
\usepackage{amsmath,amssymb}
\usepackage{booktabs}
\usepackage{longtable}
\usepackage{tabularx}
\usepackage{array}
\usepackage{enumitem}
\usepackage{xcolor}
\usepackage{hyperref}
\usepackage{caption}
\usepackage{graphicx}
\usepackage{float}

\hypersetup{
    colorlinks=true,
    linkcolor=blue,
    urlcolor=blue,
    citecolor=blue
}

\newcommand{\DevVCell}{\textsc{DevVCell}}
\newcommand{\RDEG}{\textsc{RDEG}}
\newcommand{\OT}{\textsc{OT}}
\newcommand{\GRN}{\textsc{GRN}}
\newcommand{\MVP}{\textsc{MVP}}

\title{\textbf{发育虚拟细胞 DevVCell-lite：复用 RDEG 数据资产的正式工作计划与初次执行报告}\\
\large 从阶段重建转向细胞级刺激响应、反馈恢复和发育脆弱性图谱}
\author{}
\date{\today}

\begin{document}

\maketitle

\begin{abstract}
本文档围绕 \DevVCell{} 的科学动机、原型分析和最新 \OT{}+\GRN{} 解释性补充展开。原 \RDEG{} 工作已围绕多阶段小鼠胚胎单细胞图谱完成节点构建、相邻阶段最优传输、\GRN{} 动力学、扰动模拟和 rollout 验证。\DevVCell{} 不把 \RDEG{} 作为简单扩展，而是提出一个新的计算对象：可初始化、可刺激、可反馈、可连续推演的发育虚拟细胞。当前原型先复用 \RDEG{} 的已有结果，建立阶段脆弱性、扰动优先级、反馈成本和虚拟细胞响应 proxy 等指标。第一次分析共纳入 16 个 Theiler stages 和 49 个 TF 候选，最高脆弱性阶段为 Theiler stage 15，TS14--TS19 窗口的平均脆弱性分数约为窗口外的 1.38 倍，最高扰动优先级候选为 Satb2。进一步的 \OT{}+\GRN{} 补充分析表明，调控作用可以落到 erythroid、myoblast、primitive erythroid、neural 和 cardiac muscle 等具体系统，从而把 TF--target 边与可解释的发育过程联系起来。
\end{abstract}

\section{现有数据与方法基础}

本次补充建立在既有 \RDEG{} 结果之上，而不是另起一个预测框架。现有分析已经完成多阶段小鼠胚胎单细胞图谱的系统选择、节点构建、相邻阶段 \OT{}、\GRN{} 动力学、TF 扰动、rescue、rollout 和矩阵回写。它提供了三类可以继续解释的证据：第一，相邻 Theiler stages 之间的低代价群体重排；第二，能够解释节点表达更新的 TF--target 调控边；第三，扰动、命运信息量和阶段敏感性等后验读出。

因此，本报告的重点不是说明数据和文件如何组织，而是回答一个更接近生物学的问题：在低代价群体重排和局部调控动力学耦合的过程中，哪些转录因子的调控作用最可能影响特定发育系统、细胞类型和生物过程。

\section{研究动机}

现有 \RDEG{} 工作可以回答：

\begin{quote}
胚胎阶段之间的细胞群体如何通过低代价路径完成重排，\GRN{} 动力学如何解释这种重排。
\end{quote}

\DevVCell{} 要回答的问题不同：

\begin{quote}
给定一个具体细胞状态，如果它在某个 Theiler stage 受到基因、通路或环境刺激，它会如何偏离正常轨迹，是否能通过反馈恢复，又是否会进入替代命运。
\end{quote}

因此，新项目的核心对象不是阶段节点或群体流，而是一个可初始化、可刺激、可推演的虚拟细胞 agent。这个 agent 的长期目标是输出细胞表达矩阵、命运概率、恢复成本、发育脆弱性和实验优先级。

\section{中心假说}

本文采用以下中心假说：

\begin{quote}
多阶段胚胎单细胞图谱不仅可以用于重建阶段之间的群体流向，还可以训练一个细胞级虚拟模拟器。该模拟器把正常发育视为低代价传输路径，把刺激视为对细胞状态的外部偏移，把反馈视为 \GRN{} 与发育流形共同施加的恢复力；因此，它可以预测单个细胞在不同阶段、不同扰动下的响应、恢复和命运偏移。
\end{quote}

这个假说拆成三个可执行层次：

\begin{enumerate}[label=\arabic*.]
    \item 正常发育不是随机跳变，而是沿相邻阶段之间的低代价路径演化。
    \item 刺激会把细胞推离正常发育流形，且不同阶段的可偏离程度不同。
    \item 反馈恢复不是简单回到原点，而是在 \GRN{} 约束、发育流形和刺激压力之间寻找一个低成本可行状态。
\end{enumerate}

\section{DevVCell-lite 的最小可行路线}

完整 \DevVCell{} 将来需要训练 cell encoder、transition operator、stimulus module、\GRN{} feedback head 和 decoder。当前 \DevVCell{}-lite 先不直接训练深度模型，而是执行一个低风险原型：

\begin{enumerate}[label=\arabic*.]
    \item 复用 \RDEG{} 已经生成的 stage module、TF perturbation、temporal sensitivity、rollout error 和 rescue 表。
    \item 将这些表转换为 \DevVCell{} 的指标语言。
    \item 形成三个核心结果表：stage vulnerability、perturbation priority、virtual cell rollout proxy。
    \item 生成四张初始图：阶段脆弱性、扰动优先级、反馈成本-响应幅度散点图、正常发育模块步长。
    \item 在文档中明确标注当前结果是 prototype/proxy，不等同于真实扰动实验结论。
\end{enumerate}

\section{指标设计}

\subsection{阶段脆弱性}

阶段脆弱性分数定义为：

\begin{equation}
V_s =
w_1 \tilde{S}_s
+ w_2 \tilde{E}_s
+ w_3 \tilde{D}_s
+ w_4 \tilde{O}_s
+ w_5 \tilde{M}_s,
\end{equation}

其中：

\begin{itemize}
    \item $\tilde{S}_s$ 为阶段平均 TF sensitivity 的归一化值；
    \item $\tilde{E}_s$ 为多步 rollout 误差的归一化值；
    \item $\tilde{D}_s$ 为阶段模块均值到下一阶段的正常发育步长；
    \item $\tilde{O}_s$ 为节点 outlier ratio；
    \item $\tilde{M}_s$ 为 next-step MSE；
    \item 当前权重为 $0.30, 0.25, 0.20, 0.10, 0.15$。
\end{itemize}

\subsection{扰动优先级}

扰动优先级分数定义为：

\begin{equation}
P_g =
\alpha_1 \tilde{A}_g
+ \alpha_2 \tilde{L}_g
+ \alpha_3 \tilde{W}_g
+ \alpha_4 \tilde{F}_g
+ \alpha_5 \tilde{Q}_g,
\end{equation}

其中：

\begin{itemize}
    \item $\tilde{A}_g$ 为 TF knockout 的 global effect score；
    \item $\tilde{L}_g$ 为 developmental delay score；
    \item $\tilde{W}_g$ 为 TS14--TS19 窗口内 sensitivity；
    \item $\tilde{F}_g$ 为 feedback cost；
    \item $\tilde{Q}_g$ 为 log mass shift；
    \item 当前权重为 $0.35, 0.20, 0.20, 0.15, 0.10$。
\end{itemize}

反馈成本和恢复概率的 proxy 定义为：

\begin{equation}
C_{\mathrm{feedback}}(g)=1-r_g,
\end{equation}

\begin{equation}
P_{\mathrm{recovery}}(g)=\exp(-2.5 \cdot \widetilde{C}_{\mathrm{feedback}}(g)),
\end{equation}

其中 $r_g$ 是当前 rescue 表中观察到的最佳 rescue fraction。

\section{分析工作概述}

本次工作的执行重点可以概括为五步：先将已有 \RDEG{} 派生结果转化为阶段脆弱性、扰动优先级和反馈恢复 proxy；随后检查 TF--target 调控边是否具有系统特异性；再将 TF knockout、temporal sensitivity、fate mutual information 和 hub 程度合并为解释性证据；最后把结果写成可进入正文的生物学叙事。整个过程只作为原 \OT{}+\GRN{} 框架的解释层，不新增预测模型。

任务拆解见表 \ref{tab:code-plan}。

\begin{longtable}{>{\raggedright\arraybackslash}p{4.2cm} >{\raggedright\arraybackslash}p{8.5cm} >{\raggedright\arraybackslash}p{2.5cm}}
\caption{DevVCell-lite 分析任务拆解}
\label{tab:code-plan}\\
\toprule
任务 & 说明 & 状态 \\
\midrule
\endfirsthead
\toprule
任务 & 说明 & 状态 \\
\midrule
\endhead
数据基础 & 复用多阶段胚胎单细胞图谱和既有 \RDEG{} 派生结果。 & 已完成 \\
\OT{} 结果读出 & 汇总阶段转移、rollout 误差和发育窗口敏感性。 & 已完成 \\
\GRN{} 结果读出 & 汇总 TF--target 边、系统特异边和 hub 程度。 & 已完成 \\
扰动证据整合 & 汇总 TF knockout、rescue、temporal sensitivity 和 fate mutual information。 & 已完成 \\
阶段脆弱性 & 汇总 sensitivity、rollout error、module step、outlier 和 next-step MSE。 & 已完成 \\
扰动优先级 & 汇总 knockout effect、delay、window sensitivity、feedback cost 和 mass shift。 & 已完成 \\
rollout proxy & 用阶段模块步长、阶段脆弱性和 top TF 响应构造虚拟细胞响应 proxy。 & 已完成 \\
系统发育影响 & 解释 TF 调控边主要落在哪些系统、细胞类型和生物过程。 & 已完成 \\
正式 cell-level 模型 & 接入 H5AD 子集、encoder、OT pseudo target、transition、stimulus 和 feedback head。 & 下一阶段 \\
\bottomrule
\end{longtable}

\section{第一次执行结果}

\subsection{执行概况}

第一次原型分析已经完成，摘要如下：

\begin{itemize}
    \item 纳入阶段数：16；
    \item 纳入 TF 候选数：49；
    \item 最高脆弱性阶段：Theiler stage 15；
    \item TS14--TS19 平均脆弱性分数：0.5051；
    \item 窗口外平均脆弱性分数：0.3666；
    \item TS14--TS19 与窗口外比值：1.3777；
    \item 最高扰动优先级候选：Satb2；
    \item Satb2 的 DevVCell priority score：0.8486；
    \item Satb2 的最佳 rescue TF：Zbtb16；
    \item Satb2 的 best rescue fraction：0.1214。
\end{itemize}

\subsection{阶段脆弱性初始结果}

表 \ref{tab:stage-results} 列出当前 prototype 分数最高的阶段。需要注意，Theiler stage 12 也较高，主要来自较大的模块步长和 next-step MSE；而 TS14--TS19 作为窗口整体平均值高于窗口外，支持后续把该窗口作为发育敏感窗口继续检查。

\begin{table}[H]
\centering
\caption{DevVCell-lite 阶段脆弱性初始结果}
\label{tab:stage-results}
\begin{tabularx}{\textwidth}{>{\raggedright\arraybackslash}p{3.2cm} >{\raggedleft\arraybackslash}p{2.1cm} >{\raggedright\arraybackslash}p{2.4cm} >{\raggedleft\arraybackslash}p{2.6cm} X}
\toprule
阶段 & 分数 & 窗口 & rollout 误差 & 解释 \\
\midrule
Theiler stage 15 & 0.6049 & TS14--TS19 & 0.000323 & 当前最高脆弱性阶段，处于重点敏感窗口。 \\
Theiler stage 12 & 0.6012 & outside & 0.000174 & 早期模块步长和 next-step MSE 较高，需要单独解释。 \\
Theiler stage 14 & 0.5871 & TS14--TS19 & 0.000337 & 处于窗口起点，rollout 误差较高。 \\
Theiler stage 16 & 0.5211 & TS14--TS19 & 0.000320 & 与 TS15 连续形成高脆弱区段。 \\
Theiler stage 19 & 0.5167 & TS14--TS19 & 0.000261 & 窗口后段仍保持较高分数。 \\
\bottomrule
\end{tabularx}
\end{table}

\subsection{扰动优先级初始结果}

表 \ref{tab:tf-results} 列出当前扰动优先级前五名。该结果用于后续模型训练和实验设计优先级，不直接等同于真实 knockout 实验结论。

\begin{table}[H]
\centering
\caption{DevVCell-lite 扰动优先级初始结果}
\label{tab:tf-results}
\begin{tabularx}{\textwidth}{>{\raggedright\arraybackslash}p{2.2cm} >{\raggedleft\arraybackslash}p{2.5cm} >{\raggedleft\arraybackslash}p{2.5cm} >{\raggedright\arraybackslash}p{3.2cm} X}
\toprule
TF & priority & feedback cost & peak stage & 当前解释 \\
\midrule
Satb2 & 0.8486 & 0.8786 & Theiler stage 23 & 全局效应和质量偏移最高，已有 Zbtb16 rescue proxy。 \\
Rfx4 & 0.7101 & 0.7908 & Theiler stage 14 & peak stage 落在 TS14，适合敏感窗口案例。 \\
Eya1 & 0.6875 & 0.9661 & Theiler stage 25 & rescue fraction 低，反馈成本高。 \\
Lef1 & 0.6670 & 1.0000 & Theiler stage 12 & 与 Wnt/neural 叙事相容，但当前 rescue 未观察到。 \\
Prox1 & 0.6226 & 1.0000 & Theiler stage 17 & peak stage 位于 TS14--TS19 窗口内。 \\
\bottomrule
\end{tabularx}
\end{table}

\subsection{图件结果}

\begin{figure}[H]
\centering
\includegraphics[width=0.92\textwidth]{../results/figures/stage_vulnerability.png}
\caption{DevVCell-lite 阶段脆弱性 prototype。红色为 TS14--TS19 窗口，蓝色为窗口外阶段。}
\label{fig:stage-vulnerability}
\end{figure}

\begin{figure}[H]
\centering
\includegraphics[width=0.82\textwidth]{../results/figures/perturbation_priority_top12.png}
\caption{DevVCell-lite 扰动优先级前 12 名。}
\label{fig:perturbation-priority}
\end{figure}

\section{补充分析：用原 OT+GRN 解释调控因子影响哪些系统和细胞发育}

\subsection{分析目的}

根据新的讨论，本项目不需要在这里新增预测方法。更合适的做法是直接利用原 \RDEG{} 的 \OT{}+\GRN{} 结果回答一个解释性问题：

\begin{quote}
既然模型已经学习到 TF--target 调控边，那么这些调控作用在协调演化过程中是否集中影响了某些发育系统、某些细胞类型或某些具体发育生物过程？
\end{quote}

因此，本节把原有结果表重新组织为一个 ``发育影响解释层''。它不是新模型，也不训练新的 predictor，而是把以下已有证据合并解读：

\begin{enumerate}[label=\arabic*.]
    \item \textbf{OT 约束下的阶段转移：} 相邻阶段的低代价传输给出协调演化背景。
    \item \textbf{GRN 学到的 TF--target 边：} 给出直接调控靶基因和边权。
    \item \textbf{系统特异调控边：} 给出 TF--target 边在哪些系统或细胞类型中更突出。
    \item \textbf{TF knockout 结果：} 给出全局效应、发育延迟和质量偏移。
    \item \textbf{TF--fate 信息量：} 给出 TF 与预测命运之间的信息关联。
\end{enumerate}

换言之，本节回答的不是 ``能不能用一个新方法预测细胞''，而是 ``原 OT+\GRN{} 模型已经预测和解释的发育过程，具体落在了哪些系统、细胞和调控模块上''。

\subsection{实现方式}

该补充只读取原有 \RDEG{} 输出，不读取完整 H5AD，不训练新参数。为了便于排序，我们计算了一个解释性证据汇总分数：

\begin{equation}
E_{\mathrm{OT+GRN}} =
0.35 S_{\mathrm{system-edge}}
+0.20 S_{\mathrm{KO}}
+0.15 S_{\mathrm{temporal}}
+0.15 S_{\mathrm{fateMI}}
+0.10 S_{\mathrm{hub}}
+0.05 S_{\mathrm{GRN-weight}}.
\end{equation}

这个分数只是把已有 \OT{}+\GRN{} 证据合并排序的解释性指标，不代表新的预测模型。它的作用是帮助读者看清：一个 TF 的调控边是否在某个系统中集中、这个 TF 的 knockout 是否有全局影响、它是否与命运预测有关、它是否是网络 hub。

\subsection{系统层面的结果}

表 \ref{tab:ot-grn-system-impact} 显示，原 \OT{}+\GRN{} 结果中系统特异调控边最集中的是 erythroid progenitor cell、myoblast、primitive erythroid progenitor、neural cell 和 cardiac muscle cell。这说明原方法不仅能给出阶段重排和 TF-target 边，还能把调控作用落到具体发育系统。

\begin{table}[H]
\centering
\caption{原 OT+GRN 结果读出的系统/细胞发育影响}
\label{tab:ot-grn-system-impact}
\begin{tabularx}{\textwidth}{>{\raggedright\arraybackslash}p{4.1cm} >{\raggedleft\arraybackslash}p{1.4cm} >{\raggedleft\arraybackslash}p{2cm} X}
\toprule
系统或细胞类型 & TF 数 & 边数 & 主要 TF \\
\midrule
erythroid progenitor cell & 20 & 128 & Eya1; Lef1; Prdm16; Eya4; Satb2; Pax3; Nrg1; Dach2 \\
myoblast & 20 & 128 & Eya1; Dach2; Lef1; Satb2; Satb1; Eya4; Ebf2; Pax3 \\
primitive erythroid progenitor & 19 & 123 & Eya1; Zbtb16; Lef1; Pax3; Satb1; Dach2; Satb2; Nrg1 \\
neural cell & 20 & 119 & Eya1; Satb2; Dach2; Prdm16; Eya4; Lef1; Pax3; Satb1 \\
cardiac muscle cell & 20 & 112 & Satb2; Eya1; Pax3; Eya4; Prdm16; Satb1; Dach2; Lef1 \\
myotube & 20 & 118 & Satb2; Eya1; Dach2; Lef1; Eya4; Prdm16; Satb1; Bcl11b \\
muscle precursor cell & 20 & 114 & Eya1; Dach2; Satb2; Eya4; Zbtb16; Pax3; Prdm16; Lef1 \\
hematopoietic stem cell & 20 & 101 & Satb2; Eya1; Dach2; Lef1; Prdm16; Eya4; Satb1; Pax3 \\
\bottomrule
\end{tabularx}
\end{table}

\subsection{TF 层面的结果}

表 \ref{tab:ot-grn-tf-impact} 显示，Eya1、Lef1、Prdm16、Satb2、Zbtb16、Dach2、Pax3、Satb1 和 Eya4 是当前最适合写成生物学解释案例的 TF。原因不是它们在单一表里排名靠前，而是它们同时具有若干层证据：系统特异调控边集中、\GRN{} 直接靶基因明确、knockout 或 temporal sensitivity 有响应、部分 TF 还与 fate mutual information 显著相关。

\begin{table}[H]
\centering
\caption{原 OT+GRN 证据支持的 TF--系统影响摘要}
\label{tab:ot-grn-tf-impact}
\begin{tabularx}{\textwidth}{>{\raggedright\arraybackslash}p{1.9cm} >{\raggedleft\arraybackslash}p{1.3cm} >{\raggedright\arraybackslash}p{4.2cm} >{\raggedleft\arraybackslash}p{2.1cm} X}
\toprule
TF & 系统数 & 最强系统/细胞类型 & 证据分数 & 生物学解释 \\
\midrule
Eya1 & 8 & erythroid progenitor cell & 0.7566 & 同时影响 erythroid、primitive erythroid、neural 和 myoblast，像是跨系统协调因子。 \\
Lef1 & 8 & erythroid progenitor cell & 0.6439 & 与 Wnt/FGF 叙事一致，同时在多个系统中有特异边，可能连接信号通路与命运转换。 \\
Prdm16 & 7 & erythroid progenitor cell & 0.6289 & 在 erythroid 和 neural/muscle 相关系统中均有边，可能连接代谢/应激和发育状态转换。 \\
Satb2 & 8 & cardiac muscle cell & 0.5974 & 系统边覆盖 cardiac muscle、myotube、neural 等，适合解释结构重塑和分化成熟。 \\
Zbtb16 & 8 & primitive erythroid progenitor & 0.5827 & 在 primitive erythroid 和 neural 中有系统边，可能与早期命运稳定和分支选择有关。 \\
Dach2 & 8 & myoblast & 0.5745 & 在 myoblast、neural 和 erythroid 中均突出，适合解释肌生成与神经/胚层模块耦合。 \\
Pax3 & 8 & primitive erythroid progenitor & 0.5513 & \GRN{} 中直接调控 Nrg1/Lef1 等节点，系统层面连接 primitive erythroid、cardiac 和 muscle precursor。 \\
Satb1 & 8 & primitive erythroid progenitor & 0.5396 & 与 primitive erythroid、myoblast 和 neural 的系统边共同支持其分化状态调节作用。 \\
Eya4 & 8 & erythroid progenitor cell & 0.5347 & 与 Eya1 类似，可能处于 Eya family 相关的跨系统调控模块。 \\
\bottomrule
\end{tabularx}
\end{table}

\subsection{可写进论文的生物过程叙事}

这部分结果可以把原方法的生物学解释从 ``TF knockout 有效'' 提升到 ``TF 通过哪些调控边影响哪些系统发育''。建议写成三条主线。

\paragraph{第一，造血/红系发育是最强系统信号之一。}
在系统汇总中，erythroid progenitor cell 和 primitive erythroid progenitor 位列前列，主要 TF 包括 Eya1、Lef1、Prdm16、Zbtb16、Pax3、Satb1、Satb2 和 Nrg1。这说明原 OT+\GRN{} 框架不仅在 neural MVP 上工作，也能读出红系/造血状态中较强的调控集中性。生物学上可以解释为：在胚胎协调演化过程中，红系细胞状态具有明显的命运推进和调控稳定需求，因此 TF-target 网络的系统特异边更集中。

\paragraph{第二，肌生成与心肌/骨骼肌成熟构成第二条主线。}
myoblast、myotube、muscle precursor cell 和 cardiac muscle cell 均有大量系统特异边。Satb2、Dach2、Eya1、Lef1、Pax3、Eya4 和 Prdm16 在这些系统中反复出现。结合 target 例子中的 \textit{Myl2}、\textit{Tbx15}、\textit{Adam12}、\textit{Smoc1}、\textit{Frem1}、\textit{Svep1} 等，可以把这条主线解释为 mesoderm-to-muscle specification、myogenic commitment、muscle fiber maturation 和 cardiac muscle differentiation 的调控影响。

\paragraph{第三，神经系统不是单独孤立的一条线，而是与红系和肌系共享部分调控因子。}
neural cell 中突出的 TF 包括 Eya1、Satb2、Dach2、Prdm16、Eya4、Lef1、Pax3 和 Satb1。结合 \GRN{} 中 Lef1--Fgf15、Pax3--Nrg1、Bcl11b--Cntnap2、Rorb--Epha5、Otx2--Cdh8 等直接边，可以把 neural 结果解释为 Wnt/FGF signaling、neural patterning、cell adhesion/axon guidance 和 neuronal differentiation。更重要的是，这些 TF 并不只在 neural 系统中出现，而是跨 erythroid、muscle 和 cardiac 系统出现，支持 ``协调演化'' 中共享调控模块的概念。

\subsection{对可验证 TF 的解释价值}

之前讨论的可验证 TF 不应只写成 ``敲掉以后有全局效应''。更好的表述是：

\begin{quote}
这些 TF 的扰动之所以值得验证，是因为它们在 OT 约束的发育转移背景下，具有明确的 \GRN{} 靶基因、系统特异调控边和命运信息量信号；因此可以提出具体的系统和细胞类型读出。
\end{quote}

例如：

\begin{itemize}
    \item \textbf{Lef1：} 可在 neural/Wnt-FGF 叙事中重点看 \textit{Fgf15}，同时检查 erythroid/myoblast 系统中是否出现跨系统响应。
    \item \textbf{Pax3：} \GRN{} 中直接连接 \textit{Nrg1} 和 \textit{Lef1}，适合检查 neural patterning、muscle precursor 和 primitive erythroid 的共同偏移。
    \item \textbf{Satb2：} 在 cardiac muscle、myotube 和 neural 中有系统边，适合观察结构重塑、分化成熟和命运稳定。
    \item \textbf{Eya1/Eya4：} 系统覆盖广，尤其在 erythroid 和 neural 中突出，适合作为跨系统协调因子的验证案例。
    \item \textbf{Dach2：} 在 myoblast 和 neural 中均突出，可连接肌生成与神经相关发育模块。
\end{itemize}

\begin{figure}[H]
\centering
\includegraphics[width=0.94\textwidth]{../results/figures/ot_grn_tf_system_heatmap.png}
\caption{原 OT+GRN 结果中的 TF--系统特异调控强度。颜色表示该 TF 在对应系统/细胞类型中的系统特异边绝对权重总和。}
\label{fig:ot-grn-system-heatmap}
\end{figure}

\begin{figure}[H]
\centering
\includegraphics[width=0.82\textwidth]{../results/figures/ot_grn_top_tf_developmental_impact.png}
\caption{基于原 OT+GRN 证据汇总的 TF 发育影响排序。该图仅为解释性证据排序，不是新预测模型。}
\label{fig:ot-grn-tf-impact}
\end{figure}

\section{从原型到正式模型的逐步执行计划}

\subsection{阶段一：数据子集与基因集合}

\begin{enumerate}[label=\arabic*.]
    \item 从原始 H5AD 导出 neural、mesodermal、reproductive 三个系统的小型 AnnData 子集。
    \item 每个系统保留 TS12--TS23 或 TS14--TS20 主窗口。
    \item 基因集合保留 HVG、TF、marker、pathway genes，规模控制在 2000--5000 个基因。
    \item 固定每个子集的细胞数、阶段数、donor、cell type 和基因列表，保证后续比较口径一致。
\end{enumerate}

\subsection{阶段二：正常发育虚拟细胞}

\begin{enumerate}[label=\arabic*.]
    \item 训练或加载 cell encoder，将表达矩阵映射到 latent state。
    \item 对相邻 stage 计算 mini-batch \OT{} coupling。
    \item 使用 barycentric projection 生成伪下一步目标。
    \item 训练 transition operator $F_\theta(x_i^t,t,c_i)$。
    \item rollout 1--5 步并解码回表达矩阵。
    \item 使用 heldout stage、heldout cell type 和 heldout transition 评估正常发育生成能力。
\end{enumerate}

\subsection{阶段三：刺激输入和反馈恢复}

\begin{enumerate}[label=\arabic*.]
    \item 实现 TF knockdown、TF overexpression、pathway activation/inhibition 三类 stimulus embedding。
    \item 将 \RDEG{} 的 GRN 或外部 prior 转成 feedback head 的稀疏约束。
    \item 用扰动条件下的多步 rollout 计算 response amplitude、feedback cost、fate displacement 和 recovery probability。
    \item 设计 no OT、no GRN、no feedback、no stimulus embedding 等消融实验。
\end{enumerate}

\subsection{阶段四：外部扰动验证}

\begin{enumerate}[label=\arabic*.]
    \item 接入公开 perturbation 数据，优先考虑和 neural、Wnt、RA、FGF、PGC-like 或 differentiation 相关的数据。
    \item 与 CellOT、CPA、scGPT/Geneformer embedding + MLP 进行比较。
    \item 指标包括 DEG overlap、logFC Spearman、Wasserstein distance、cell-state transition accuracy 和 pathway recovery。
    \item 若条件允许，设计 2--3 个小规模实验验证：RA/Wnt 敏感窗口、Lef1 perturbation、Glis3/Tbx4 rescue。
\end{enumerate}

\section{风险控制}

\begin{enumerate}[label=\arabic*.]
    \item \textbf{真实扰动数据不足：} 当前结果必须标注为 counterfactual simulation 或 proxy，不能写成已验证事实。
    \item \textbf{虚拟细胞概念过大：} 论文中应限定为 embryo-context developmental virtual cell，而不是 universal virtual cell。
    \item \textbf{与 RDEG 重叠：} 第一篇讲阶段重建和群体流，第二篇讲单细胞刺激响应和反馈恢复。
    \item \textbf{模型过重：} 先完成 DevVCell-lite 和轻量 cell-level model，再评估是否引入 scGPT、Geneformer 或更大模型 embedding。
\end{enumerate}

\section{当前结论}

本次结果支持将 TS14--TS19 作为重点敏感窗口继续推进，同时提示 Theiler stage 12 也需要作为早期高变化阶段单独建模。扰动优先级上，Satb2、Rfx4、Eya1、Lef1 和 Prox1 是当前最值得进入下一步训练或验证的候选。更重要的是，原 \OT{}+\GRN{} 结果已经能够进一步解释 TF 调控作用的系统落点：Eya1、Lef1、Prdm16、Satb2、Dach2、Pax3 等 TF 的系统特异调控边分别指向 erythroid、myoblast、primitive erythroid、neural 和 cardiac muscle 等发育过程。下一阶段应在保持原 \OT{}+\GRN{} 框架的基础上，把这些系统级解释转化为更具体的细胞类型读出和实验验证设计。

\begin{thebibliography}{99}

\bibitem{rdeg_validation}
本项目内部报告：\textit{RDEG 宏观--微观联合推演框架的可验证发现汇总：GRN 新边、谱系推演、基因扰动与滚动预测误差窗口}，2026。

\bibitem{rdeg_integrated}
本项目内部报告：\textit{单细胞胚胎发育图谱中的低代价群体重排、基因调控动力学与严格谱系证据边界：从宏观发育重建到 PGC--Leydig 生殖系统案例的统一分析}，2026。

\bibitem{waddingtonot}
Schiebinger, G. et al.
Optimal-transport analysis of single-cell gene expression identifies developmental trajectories in reprogramming.
\textit{Cell}, 2019.

\bibitem{cuturi}
Cuturi, M.
Sinkhorn distances: Lightspeed computation of optimal transport.
\textit{Advances in Neural Information Processing Systems}, 2013.

\bibitem{cellot}
Bunne, C. et al.
Learning single-cell perturbation responses using neural optimal transport.
\textit{Nature Methods}, 2023.

\bibitem{cpa}
Lotfollahi, M. et al.
Predicting cellular responses to complex perturbations in high-throughput screens.
\textit{Molecular Systems Biology}, 2023.

\bibitem{celloracle}
Kamimoto, K. et al.
Dissecting cell identity via network inference and in silico gene perturbation.
\textit{Nature}, 2023.

\bibitem{geneformer}
Theodoris, C. V. et al.
Transfer learning enables predictions in network biology.
\textit{Nature}, 2023.

\bibitem{scgpt}
Cui, H. et al.
scGPT: toward building a foundation model for single-cell multi-omics using generative AI.
\textit{Nature Methods}, 2024.

\end{thebibliography}

\end{document}
